% =====================================================================
%  第 1 节  引言（原书 p.2–3）
% =====================================================================
\section{引言}\label{sec:1}

% [Tutor Note: 本节是全文的路标，一句话概括：把 Fukaya 范畴在余切丛 $T^*X$ 上取半经典极限 $\epsilon\to0$，全纯圆盘的计数就退化为梯度流树的计数；本文补上“对象到自身”的态射，并给出 $A_\infty$ 关系式的两个非全纯证明。]

\begin{takeaway}
\begin{itemize}
  \item \textbf{物理图像}：Fukaya 范畴用全纯圆盘的计数定义 Lagrange 子流形之间的高阶复合。在余切丛 $T^*X$ 上取 $L_a=\mathrm{graph}(\epsilon\,\mathrm{d}F_a)$ 并令 $\epsilon\to0$，全纯圆盘投影到 $X$ 后收缩为梯度流树，全纯计数因而退化为 Morse 理论中的计数。此即 Fukaya--Morse $\Ainf$ 范畴。
  \item \textbf{数学工具}：$\Ainf$ 预范畴、Morse 链复形、生成函数、带状树、横截性。
  \item \textbf{主要结论}：该构造原本缺“对象到自身的态射”。本文取 $\Mor(F_a,F_a)=C^{\mathrm{sing}}_\bullet(X;\mathbb Z)$ 补上这一环，得到增强的 Fukaya--Morse $\Ainf$ 范畴，并给出 $\Ainf$ 关系式的两种证明：BF 理论的有效作用量（二次量子化）与度量树上的高阶拓扑量子力学（一次量子化）。两个证明都不使用全纯圆盘。
\end{itemize}
\end{takeaway}

Fukaya 在 \cite{fukaya} 中引入了一个 $\Ainf$ 范畴：对象是辛流形中的 Lagrange 子流形，态射是这些子流形交点的形式线性组合，高阶复合则以这些子流形为边界的全纯圆盘\footnote{这里沿用一个习惯的术语滥用：下文说“全纯”，指的其实是“伪全纯”，即相对于辛流形上一个相容的殆复结构而言。}的计数来定义。\footnote{细节参看 \cite{fukayaoh,seidel}；本主题的入门导引见 \cite{auroux}。}若辛流形是余切丛 $T^*X$，且 Lagrange 子流形都能投影到底流形 $X$ 上，则 $N$ 个不同的 Lagrange 子流形 $L_a$（假定它们都是正合的）可用 $N$ 个生成函数 $F_a$ 来描述，$L_a$ 由方程
\begin{equation*}
  p_i=\frac{\partial F_a}{\partial X^i}
\end{equation*}
给出，其中 $p_i$ 是与底流形坐标 $X^i$ 对偶的纤维坐标。于是 Lagrange 子流形的交点投影为生成函数之差的临界点 $P_{\alpha,ab}$：
\begin{equation*}
  \mathrm{d}(F_a-F_b)=0.
\end{equation*}
若作重标度 $F_a\to\epsilon F_a$，则全纯圆盘到底流形的投影趋于由生成函数之差的梯度轨道拼成的带状树。计数这类带状树就确定了某个 $\Ainf$ 范畴的复合映射。为与辛流形中由全纯圆盘计数给出的“通常的”Fukaya 范畴相区别，我们称这一结构为“Fukaya--Morse $\Ainf$ 范畴”。细节见原始论文 \cite{fukaya,fukayaoh}，以及 \cite{kontsevichsoibelman}。

% [Tutor Note: “重标度 $F_a\to\epsilon F_a$”是全篇第一个易错点：$\epsilon\to0$ 时 $L_a=\mathrm{graph}(\epsilon\,\mathrm{d}F_a)$ 退化到零截面，交点 $P_{\alpha,ab}$ 趋于 $\mathrm{d}(F_a-F_b)=0$ 的临界点，全纯圆盘的半径随之趋于零。作者在此省略了“一列全纯圆盘收敛到梯度流树”的紧性论证，标准做法是用 Gromov 紧性与 $T^*X$ 上近全纯曲线的局部描述，见 Fukaya--Oh \cite{fukayaoh} 与 Seidel \cite{seidel} 第 9 章。]

这一构造所缺的，是对象到自身态射的描述。我们想补上这一空缺，给出一个包含这类态射的 $\Ainf$ 结构。本文将给出构造，并给出 $\Ainf$ 关系式的两个不同证明。两个证明都不涉及全纯圆盘。

% [Tutor Note: “两个证明都不涉及全纯圆盘”是全文纲领，值得停一下。其一只用 BV 形式的纤维积分与树形求和（等价于 Kontsevich--Soibelman 同伦转移定理，见 \cite{kontsevichsoibelman}）；其二只用度量树模空间上一个闭微分形式与 Stokes 定理。换言之，作者把“辛几何”这一层完全剥离，只留下“生成函数之差是 Morse--Smale 函数”与“梯度流”两件事。这也解释了为什么证明能同时给出整系数版本与一般的环系数版本。]

\subsection{结果概述}\label{sec:1.1}

\begin{takeaway}
\begin{itemize}
  \item \textbf{物理图像}：三点新意分别对应“补上自同态”（第~\ref{sec:2} 节）、“把整个结构翻译成 BF 理论的有效作用量”（第~\ref{sec:3} 节）、“把同一批数据翻译成树上的高阶拓扑量子力学”（第~\ref{sec:4} 节）。
  \item \textbf{数学工具}：奇异链与余维数分次、Batalin--Vilkovisky（BV）主方程、模空间上的积分与周期的 Stokes 定理。
  \item \textbf{主要结论}：$\Ainf$ 关系式在这里可看作 BV 主方程在“慢场”上的投影，也可看作“度量树模空间上闭形式在边界处分解”的 Stokes 推论。
\end{itemize}
\end{takeaway}

与已有文献相比，本文的新意有以下三点。

\begin{enumerate}
  \item[(A)] \textbf{第~\ref{sec:2} 节}：用一个对象自同态的模型来增强 Fukaya--Morse $\Ainf$ 范畴，即取
  \begin{equation*}
    \Mor(F_a,F_a)=C^{\mathrm{sing}}_\bullet(X,\mathbb Z),
  \end{equation*}
  也就是 $X$ 上的奇异链。高阶复合由一个计数几何问题定义：在 $X$ 中计数由梯度轨道构成的树，这些树连接给定的一组临界点，并经过给定的一组链。
  \item[(B)] \textbf{第~\ref{sec:3} 节}：用拓扑规范理论（BF 理论）在特定规范下算出的有效作用量，来描述增强 Fukaya--Morse $\Ainf$ 范畴的数据。在此框架中，$\Ainf$ 范畴的结构关系式由有效作用量的 Batalin--Vilkovisky 主方程推出。
  \item[(C)] \textbf{第~\ref{sec:4} 节}：用高阶拓扑量子力学来描述同一批数据，即给出度量树模空间上的一个微分形式 $I$。用这套语言，$\Ainf$ 复合映射是形式 $I$ 的周期，$\Ainf$ 关系式则由模空间上的 Stokes 定理与 $I$ 在模空间边界处的分解性质推出。此证明是上同调场论中 WDVV 方程证明的一维类比，见 \cite{kontsevichmanin}。
\end{enumerate}

% [Tutor Note: (B) 的逻辑方向容易说反。BV 主方程 $\{S_{\mathrm{eff}},S_{\mathrm{eff}}\}=0$ 直接给出的是 $L_\infty$ 关系式（斜对称化之后的结果）；要得到 $A_\infty$ 关系式，必须把系数代数换成上三角矩阵代数 $\mathcal A$ 再张量，见注 3.2 与式 (47)。作者在 (B) 里没有点破这一层，读者若不留意会误以为 $A_\infty$ 关系式可以直接从 BV 主方程读出。]

% [Tutor Note: (A) 中“取奇异链”不是任意选择：$F_a-F_a\equiv0$ 的临界点集是整个 $X$，故自同态的 Morse 复数形式上就是 $X$ 上的链；取奇异链（而非只取点）使复合映射 (4) 能作用在“被看作分布”的链上。若改取 de Rham 形式或上链，会付出别的代价，见第 2.7 节与我在此处的批注。]
\endinput
